% c03 原书 p.19–29（PDF p.026–036）
% 末页延至原书 p.30 顶部：§1.2 的 Exercises 整块（1–7 题）按「条目起始页」规则归本块，已整块译完。
\subsection{概形与态射}
\label{sec:1.2}

全书假定：所有环都是交换环且含单位元 $1$，所有环同态都把 $1$ 送到 $1$。

对每个环 $A$，令 $\Spec A$ 为 $A$ 的\emph{素理想}\idx{sulixiang}{素理想}构成的集合。对 $A$ 的每个理想 $\fa$，定义
\[
V(\fa)=\{\fp\in \Spec A\mid \fa\subset \fp\}.
\]
下面的命题留给读者证明：
\begin{proposition}\label{it:1.2.1}
\begin{enumerate}[label=(\roman*)]
  \item $V(0)=\Spec A$，$V(A)=\emptyset$。
  \item 若 $\fa\subset \fb$，则 $V(\fa)\supset V(\fb)$。
  \item 对 $A$ 的任意一族理想 $\fa_i$（$i\in I$），有 $\bigcap_{i\in I}V(\fa_i)=V\bigl(\sum_{i\in I}\fa_i\bigr)$。
  \item $V(\fa)\cup V(\fb)=V(\fa\fb)=V(\fa\cap\fb)$。
\end{enumerate}
\end{proposition}

这个命题表明，$\Spec A$ 是拓扑空间，其闭集恰为理想 $\fa$ 对应的 $V(\fa)$。这个拓扑称为 $\Spec A$ 上的 \emph{Zariski 拓扑}\idx{zariskituopu}{Zariski 拓扑}。

\begin{proposition}\label{it:1.2.2}
\begin{enumerate}[label=(\roman*)]
  \item $\Spec A=\emptyset$ 当且仅当 $A$ 中 $0=1$。
  \item 对 $A$ 的任意理想 $\fa$，定义 $\fa$ 的\emph{幂零根}\idx{milinggen}{幂零根}为理想
  \[
  \sqrt{\fa}=\{a\in A\mid a^n\in\fa\ \text{对某个自然数 } n\}.
  \]
  则 $V(\fa)=V(\sqrt{\fa})$。
  \item $\sqrt{\fa}=\bigcap_{\fp\in V(\fa)}\fp$。
  \item 对 $A$ 的任意理想 $\fa$ 与 $\fb$，有 $V(\fa)\subset V(\fb)$ 当且仅当 $\sqrt{\fa}\supset\sqrt{\fb}$。
\end{enumerate}
\end{proposition}

\begin{proof}
(i) 若 $0=1$，则 $A=0$，从而 $\Spec A=\emptyset$。设 $0\ne 1$。令
\[
\mathcal{S}=\{\fa\mid \fa \text{ 是 } A \text{ 的理想，且 } 1\notin\fa\}.
\]
则 $(0)\in\mathcal{S}$，故 $\mathcal{S}$ 非空。易证：$\mathcal{S}$ 中按包含关系全序的任意子集 $\mathcal{T}$ 在 $\mathcal{S}$ 中有上界 $\bigcup_{\fa\in\mathcal{T}}\fa$。于是由 Zorn 引理，$\mathcal{S}$ 有极大元，即 $A$ 有极大理想\idx{jidalixiang}{极大理想}。极大理想都是素理想。所以 $\Spec A$ 非空。

(ii) 直接由 $V(\fa)$ 的定义得到。

(iii) 显然 $\sqrt{\fa}\subset\bigcap_{\fp\in V(\fa)}\fp$。设 $f\notin\sqrt{\fa}$。则在 $(A/\fa)_f$ 中有 $0\ne 1$。由 (i)，可取 $(A/\fa)_f$ 的一个素理想 $\fq$。设 $\fp$ 是 $\fq$ 在典范同态 $A\to(A/\fa)_f$ 下的原像。则 $\fp\in V(\fa)$ 但 $f\notin\fp$。故 $f\notin\bigcap_{\fp\in V(\fa)}\fp$。

(iv) 由 (ii)、(iii) 及 \ref{it:1.2.1}(ii) 得到。
\end{proof}

\begin{tnote}
命题 \ref{it:1.2.2}(ii) 的记号要留意：$\sqrt{\fa}$ 是理想，而花括号里的 $a$ 是环元素，两者在原书中都印成字母「a」，容易看混——元素是普通斜体 $a$，理想是 Fraktur 体 $\fa$。另外 (iii) 说明 $\Spec A$ 的闭集只依赖 $\fa$ 的幂零根，因此闭子集与「根理想」一一对应；这正是本节把几何对象（闭集）与代数对象（根理想）对应起来的出发点。(i) 则保证了空集也作为闭集出现，Zariski 拓扑才定义得干净。[依据：原书 p.19 命题 1.2.2 及其证明]
\end{tnote}

\begin{proposition}\label{it:1.2.3}
对任意 $f\in A$，子集
\[
D(f)=\Spec A-V((f))=\{\fp\in \Spec A\mid f\notin\fp\}
\]
是开集。这类开集构成 $\Spec A$ 的 Zariski 拓扑的一组基，也就是说，$\Spec A$ 的任意开集都是这类开集的并。此外，$D(f)$ 是\emph{拟紧}\idx{nijin}{拟紧}的。
\end{proposition}

\begin{proof}
由于 $D(f)$ 是闭集 $V((f))$ 的补集，它是开集。设 $U$ 是 $\Spec A$ 的开子集，$\fp\in U$。则对 $A$ 的某个理想 $\fa$ 有 $U=\Spec A-V(\fa)$。我们有 $\fa\not\subset\fp$。取 $f\in\fa-\fp$。则 $D(f)$ 是 $\fp$ 的邻域且包含在 $U$ 中。所以开集 $D(f)$（$f\in A$）构成 $\Spec A$ 的一组基。

设 $D(f)\subset\bigcup_{i\in I}D(f_i)$，其中 $f_i$（$i\in I$）是 $A$ 中的一族元素。则
\[
V((f))\supset\bigcap_{i\in I}V((f_i))=V\Bigl(\sum_{i\in I}(f_i)\Bigr).
\]
由 \ref{it:1.2.2}(iv)，有 $\sqrt{(f)}\subset\sqrt{\sum_{i\in I}(f_i)}$。于是可找到自然数 $n$，使得
\[
f^n=a_{i_1}f_{i_1}+\cdots+a_{i_k}f_{i_k},
\]
其中 $i_1,\dots,i_k\in I$，$a_{i_1},\dots,a_{i_k}\in A$。这说明 $D(f)\subset D(f_{i_1})\cup\cdots\cup D(f_{i_k})$。所以 $D(f)$ 是拟紧的。
\end{proof}

在 $\Spec A$ 上定义环层 $\cO_{\Spec A}$ 如下：对 $\Spec A$ 的每个开子集 $U$，$\cO_{\Spec A}(U)$ 由满足以下两个条件的映射 $s\colon U\to\coprod_{\fp\in\Spec A}A_{\fp}$ 组成：
\begin{enumerate}[label=(\alph*)]
  \item 对每个 $\fp\in U$，有 $s(\fp)\in A_{\fp}$。
  \item 对每个 $\fp\in U$，存在 $\fp$ 的包含在 $U$ 中的邻域 $U_{\fp}$ 及 $a,f\in A$，使得对每个 $\fq\in U_{\fp}$，有 $f\notin\fq$ 且 $s(\fq)=\frac{a}{f}\in A_{\fq}$。
\end{enumerate}

对开子集的每个包含 $V\subset U$，定义 $\cO_{\Spec A}(U)\to\cO_{\Spec A}(V)$ 为映射的限制。我们称二元组 $(\Spec A,\cO_{\Spec A})$ 为 $A$ 的\emph{谱}\idx{pu}{谱}。若无混淆，常把 $\cO_{\Spec A}$ 记为 $\cO$。

\begin{tnote}
截面 $s$ 的取值落在不交并 $\coprod_{\fp\in\Spec A}A_{\fp}$ 中，即每个 $s(\fp)$ 都身在各自的局部化 $A_{\fp}$ 里；而命题 \ref{it:1.2.4}(ii) 的证明把同一个截面写成乘积 $\prod_{\fp\in\Spec A}A_{\fp}$ 的元素。这里两者可以互换：条件 (a) 已经指定了 $s(\fp)\in A_{\fp}$，所以「一族按点取值的元素」与「不交并中的元素」是同一件事。原书前后混用这两个记号，并非笔误。[依据：原书 p.20 结构层的定义与 p.21 命题 1.2.4 的证明]
\end{tnote}

\begin{proposition}\label{it:1.2.4}
\begin{enumerate}[label=(\roman*)]
  \item 对每个 $\fp\in\Spec A$，有典范同构 $\cO_{\fp}\cong A_{\fp}$。
  \item 对每个 $f\in A$，有典范同构 $\cO(D(f))\cong A_f$。特别地，取 $f=1$ 得 $\cO(\Spec A)\cong A$。
\end{enumerate}
\end{proposition}

\begin{proof}
(i) 对 $\fp$ 的每个邻域 $U$，有同态 $\cO(U)\to A_{\fp}$，定义为 $s\mapsto s(\fp)$。这些同态诱导同态
\[
\cO_{\fp}=\dlim_{\fp\in U}\cO(U)\to A_{\fp}.
\]
容易证明它是满射。为证它是单射，设 $(s,U)\in\cO(U)$ 满足 $s(\fp)=0$。取 $\fp$ 的包含在 $U$ 中的邻域 $U_{\fp}$ 及 $a,f\in A$，使对每个 $\fq\in U_{\fp}$ 有 $f\notin\fq$ 且 $s(\fq)=\frac{a}{f}\in A_{\fq}$。由 $s(\fp)=0$ 得 $\frac{a}{f}=0$ 在 $A_{\fp}$ 中成立。故存在 $t\notin\fp$ 使 $ta=0$。于是对每个 $\fq\in D(f)\cap D(t)$，在 $A_{\fq}$ 中有 $\frac{a}{f}=\frac{ta}{tf}=0$。因此 $s|_{U_{\fp}\cap D(t)}=0$。注意 $U_{\fp}\cap D(t)$ 是 $\fp$ 的邻域。这说明 $\cO_{\fp}\to A_{\fp}$ 是单射。

(ii) 我们来证明同态
\[
A_f\to\cO(D(f)),
\]
它把 $\frac{a}{f^k}\in A_f$ 送到截面
\[
D(f)\to\prod_{\fp\in\Spec A}A_{\fp},\qquad \fp\mapsto \frac{a}{f^k}\in A_{\fp},
\]
是双射。

单射性：设 $\frac{a}{f^k}$ 落在上述同态的核中。则对每个 $\fp\in D(f)$，在 $A_{\fp}$ 中有 $\frac{a}{f^k}=0$。所以对每个 $\fp\in D(f)$，存在 $t\notin\fp$ 使 $ta=0$。于是
\[
\fp\in D(f)\Rightarrow \operatorname{Ann}(a)\not\subset\fp,
\]
其中 $\operatorname{Ann}(a)=\{r\in A\mid ra=0\}$ 是 $a$ 的零化子。所以
\[
\fp\notin V((f))\Rightarrow \fp\notin V(\operatorname{Ann}(a)).
\]
因此
\[
V(\operatorname{Ann}(a))\subset V((f)).
\]
由 \ref{it:1.2.2}(iv)，有 $f\in\sqrt{\operatorname{Ann}(a)}$，即对某个自然数 $n$ 有 $f^na=0$。所以 $\frac{a}{f^k}=0$ 在 $A_f$ 中。这证明了单射性。

满射性：设 $s\colon D(f)\to\prod_{\fp\in\Spec A}A_{\fp}$ 是 $\cO(D(f))$ 中的截面。可以找到 $D(f)$ 的开覆盖 $\{U_i\}_{i\in I}$ 及 $a_i,f_i\in A$，使得对每个 $\fq\in U_i$，有 $f_i\notin\fq$ 且 $s(\fq)=\frac{a_i}{f_i}\in A_{\fq}$。由 \ref{it:1.2.3}，可设 $U_i=D(g_i)$，其中 $g_i\in A$；又可设该覆盖是有限覆盖。则对任意 $\fq\in D(g_i)$ 有 $f_i\notin\fq$。故 $D(g_i)\subset D(f_i)$。由 \ref{it:1.2.2}(iv)，有 $\sqrt{(g_i)}\subset\sqrt{(f_i)}$。于是对某个自然数 $k_i$ 与某个 $h_i\in A$ 有 $g_i^{k_i}=h_if_i$。注意对每个 $\fq\in D(g_i)$ 有 $h_i,f_i\notin\fq$，且
\[
s(\fq)=\frac{a_i}{f_i}=\frac{h_ia_i}{h_if_i}=\frac{h_ia_i}{g_i^{k_i}}
\]
在 $A_{\fq}$ 中成立。此外有 $D(g_i)=D(g_i^{k_i})$。把 $g_i$ 换成 $g_i^{k_i}$、$a_i$ 换成 $h_ia_i$，可以设 $D(f)$ 有有限覆盖 $\{D(g_i)\}_{i\in I}$，并且对每个 $\fq\in D(g_i)$ 有 $s(\fq)=\frac{a_i}{g_i}\in A_{\fq}$。

对每个 $\fq\in D(g_i)\cap D(g_j)=D(g_ig_j)$，在 $A_{\fq}$ 中有 $\frac{a_i}{g_i}=\frac{a_j}{g_j}$。所以对每个 $\fq\in D(g_ig_j)$，存在 $t\notin\fq$ 使 $t(a_ig_j-a_jg_i)=0$。于是对每个 $\fq\in D(g_ig_j)$ 有 $\operatorname{Ann}(a_ig_j-a_jg_i)\not\subset\fq$；也就是说，若 $\fq\notin V((g_ig_j))$，则 $\fq\notin V(\operatorname{Ann}(a_ig_j-a_jg_i))$。因此
\[
V(\operatorname{Ann}(a_ig_j-a_jg_i))\subset V((g_ig_j)).
\]
由 \ref{it:1.2.2}(iv)，有 $g_ig_j\in\sqrt{\operatorname{Ann}(a_ig_j-a_jg_i)}$。取充分大的整数 $k$，使得对每对 $i,j$ 都有
\[
(g_ig_j)^k(a_ig_j-a_jg_i)=0.
\]
（我们的覆盖是有限覆盖。）我们有 $\frac{a_i}{g_i}=\frac{a_ig_i^k}{g_i^{k+1}}$。把 $g_i$ 换成 $g_i^{k+1}$、$a_i$ 换成 $a_ig_i^k$，可以设 $\{D(g_i)\}_{i\in I}$ 是 $D(f)$ 的有限覆盖，且对每个 $\fq\in D(g_i)$ 有 $s(\fq)=\frac{a_i}{g_i}\in A_{\fq}$，并对每对 $i,j$ 有 $a_ig_j=a_jg_i$。

由 $D(f)=\bigcup_iD(g_i)$ 得 $V((f))=\bigcap_iV((g_i))=V\bigl(\sum_i(g_i)\bigr)$。由 \ref{it:1.2.2}(iv)，存在自然数 $k$ 与 $b_i\in A$ 使 $f^k=\sum_ib_ig_i$。对每个 $j$ 有
\[
a_jf^k=a_j\sum_ib_ig_i=\sum_ib_ia_jg_i=\sum_ib_ia_ig_j=\Bigl(\sum_ib_ia_i\Bigr)g_j.
\]
所以有 $\frac{a_j}{g_j}=\frac{\sum_ib_ia_i}{f^k}$。于是按开头定义的同态 $A_f\to\cO(D(f))$，$\frac{\sum_ib_ia_i}{f^k}$ 被送到截面 $s\in\cO(D(f))$。这证明了满射性。
\end{proof}

\emph{环层空间}\idx{huancengkongjian}{环层空间}是一个二元组 $(X,\cO_X)$，由拓扑空间 $X$ 与 $X$ 上的环层 $\cO_X$ 组成。我们称 $X$ 为\emph{底拓扑空间}\idx{dituopukongjian}{底拓扑空间}，称 $\cO_X$ 为\emph{结构层}\idx{jiegouceng}{结构层}。若对每个 $P\in X$，$\cO_X$ 在 $P$ 处的茎 $\cO_{X,P}$ 是\emph{局部环}\idx{jubuhuan}{局部环}，则称 $(X,\cO_X)$ 是\emph{局部环层空间}\idx{jubuhuancengkongjian}{局部环层空间}。（回忆：环 $A$ 称为局部环，若它只有一个极大理想；这等价于说 $A$ 中不是单位的元素构成 $A$ 的一个理想。）设 $\fm_P$ 是 $\cO_{X,P}$ 的极大理想。我们称 $k(P)=\cO_{X,P}/\fm_P$ 为 $(X,\cO_X)$ 在 $P$ 处的\emph{剩余域}\idx{shengyuyu}{剩余域}。由 \ref{it:1.2.4}(i)，对任意环 $A$，谱 $(\Spec A,\cO_{\Spec A})$ 是局部环层空间。若 $X$ 是拓扑空间，$\cC_X$ 是 $X$ 上连续函数构成的层，则 $(X,\cC_X)$ 也是局部环层空间。

设 $(X,\cO_X)$ 与 $(Y,\cO_Y)$ 是两个环层空间。从 $(X,\cO_X)$ 到 $(Y,\cO_Y)$ 的\emph{态射}是一个二元组 $(f,f^\sharp)$，由连续映射 $f\colon X\to Y$ 与层态射 $f^\sharp\colon\cO_Y\to f_*\cO_X$ 组成。对每个点 $P\in X$，有同态 $(f_*\cO_X)_{f(P)}\to\cO_{X,P}$，定义为
\[
(f_*\cO_X)_{f(P)}=\dlim_{f(P)\in V}(f_*\cO_X)(V)=\dlim_{P\in f^{-1}(V)}\cO_X(f^{-1}(V))\to\dlim_{P\in U}\cO_X(U)=\cO_{X,P}.
\]
把此同态与 $f^\sharp$ 诱导的同态 $\cO_{Y,f(P)}\to(f_*\cO_X)_{f(P)}$ 复合，得到同态
\[
f^\sharp_P\colon \cO_{Y,f(P)}\to\cO_{X,P}.
\]

设 $(X,\cO_X)$ 与 $(Y,\cO_Y)$ 是两个局部环层空间。\emph{局部环层空间的态射} $(f,f^\sharp)\colon(X,\cO_X)\to(Y,\cO_Y)$ 是环层空间的态射，且对每个 $P\in X$，同态 $f^\sharp_P\colon\cO_{Y,f(P)}\to\cO_{X,P}$ 是局部同态。（回忆：局部环的\emph{局部同态}\idx{jubutongtai}{局部同态} $f\colon A\to B$ 是环同态，满足 $f^{-1}(\fm_B)=\fm_A$；等价地，$f(\fm_A)\subset\fm_B$，其中 $\fm_A$ 与 $\fm_B$ 分别是 $A$ 与 $B$ 的极大理想。）设 $(f,f^\sharp)\colon(X,\cO_X)\to(Y,\cO_Y)$ 与 $(g,g^\sharp)\colon(Y,\cO_Y)\to(Z,\cO_Z)$ 是局部环层空间的两个态射。我们如下定义它们的复合：在底拓扑空间上，态射是复合 $gf\colon X\to Z$；在结构层上，态射是复合
\[
\cO_Z\xrightarrow{g^\sharp}g_*\cO_Y\xrightarrow{g_*(f^\sharp)}g_*f_*\cO_X=(gf)_*\cO_X.
\]
局部环层空间的\emph{同构} $(f,f^\sharp)\colon(X,\cO_X)\to(Y,\cO_Y)$ 是有双边逆的态射。这等价于说，$f\colon X\to Y$ 是拓扑空间的同胚，且 $f^\sharp\colon\cO_Y\to f_*\cO_X$ 是层的同构。

与某个环 $A$ 的 $(\Spec A,\cO_{\Spec A})$ 同构的局部环层空间称为\emph{仿射概形}\idx{fangshegaixing}{仿射概形}。\emph{概形}\idx{gaixing}{概形} $(X,\cO_X)$ 是局部环层空间，且存在 $X$ 的开覆盖 $\{U_i\}_{i\in I}$，使得每个 $(U_i,\cO_X|_{U_i})$ 都是仿射概形。常把 $(X,\cO_X)$ 记为 $X$。我们把概形的态射\idx{gaixingtaishe}{概形的态射}定义为局部环层空间的态射。常把概形的态射 $(f,f^\sharp)\colon(X,\cO_X)\to(Y,\cO_Y)$ 记为 $f\colon X\to Y$。

\begin{proposition}\label{it:1.2.5}
\begin{enumerate}[label=(\roman*)]
  \item 设 $\phi\colon A\to B$ 是环同态。则 $\phi$ 典范地诱导局部环层空间的态射
  \[
  (f,f^\sharp)\colon(\Spec B,\cO_{\Spec B})\to(\Spec A,\cO_{\Spec A}).
  \]
  \item 局部环层空间的任意态射 $(f,f^\sharp)\colon(\Spec B,\cO_{\Spec B})\to(\Spec A,\cO_{\Spec A})$ 都由此得到。
\end{enumerate}
\end{proposition}

\begin{proof}
(i) 给定环同态 $\phi\colon A\to B$，定义 $f\colon\Spec B\to\Spec A$ 为
\[
f(\fq)=\phi^{-1}(\fq)
\]
对每个 $\fq\in\Spec B$。对 $A$ 的每个理想 $\fa$，有 $f^{-1}(V(\fa))=V(\fa B)$，其中 $\fa B$ 是由 $\phi(\fa)$ 生成的 $B$ 的理想。故 $f$ 连续。

对 $\Spec A$ 的每个开子集 $V$ 与 $\cO_{\Spec A}(V)$ 中的每个截面 $s\colon V\to\prod_{\fp\in\Spec A}A_{\fp}$，定义 $\cO_{\Spec B}(f^{-1}(V))$ 中的截面 $f^\sharp(s)$ 为
\[
f^\sharp(s)\colon f^{-1}(V)\to\prod_{\fq\in\Spec B}B_{\fq},\qquad (f^\sharp(s))(\fq)=\phi_{\fq}(s(f(\fq)))
\]
其中 $\fq\in f^{-1}(V)$，且 $\phi_{\fq}\colon A_{\phi^{-1}(\fq)}\to B_{\fq}$ 是由 $\phi$ 诱导的同态。这样就得到层态射 $f^\sharp\colon\cO_{\Spec A}\to f_*(\cO_{\Spec B})$。可以验证：在 \ref{it:1.2.4}(i) 给出的等同下，同态 $f^\sharp_{\fq}\colon\cO_{\Spec A,f(\fq)}\to\cO_{\Spec B,\fq}$ 与 $\phi_{\fq}$ 一致。所以 $f^\sharp_{\fq}$ 是局部同态，从而 $(f,f^\sharp)$ 是概形的态射。

(ii) 设 $(f,f^\sharp)\colon(\Spec B,\cO_{\Spec B})\to(\Spec A,\cO_{\Spec A})$ 是局部环层空间的态射。定义 $\phi\colon A\to B$，使下图交换：
\[
\begin{tikzcd}
A \arrow[r,"\phi"] \arrow[d,"\cong"'] & B \arrow[d,"\cong"] \\
\cO_{\Spec A}(\Spec A) \arrow[r,"f^\sharp"'] & (f_*\cO_{\Spec B})(\Spec A)
\end{tikzcd}
\]
其中竖直箭头由 \ref{it:1.2.4}(ii) 定义。对每个 $\fq\in\Spec B$，定义 $\phi'_{\fq}\colon A_{f(\fq)}\to B_{\fq}$，使下图交换：
\[
\begin{tikzcd}
\cO_{\Spec A,f(\fq)} \arrow[r,"f^\sharp_{\fq}"] \arrow[d,"\cong"'] & \cO_{\Spec B,\fq} \arrow[d,"\cong"] \\
A_{f(\fq)} \arrow[r,"\phi'_{\fq}"'] & B_{\fq}
\end{tikzcd}
\]
其中竖直箭头由 \ref{it:1.2.4}(i) 定义。由于 $f^\sharp_{\fq}$ 是局部同态，$\phi'_{\fq}$ 也是局部的。下图交换：
\[
\begin{tikzcd}
\cO_{\Spec A}(\Spec A) \arrow[r,"f^\sharp"] \arrow[d] & \cO_{\Spec B}(\Spec B) \arrow[d] \\
\cO_{\Spec A,f(\fq)} \arrow[r,"f^\sharp_{\fq}"'] & \cO_{\Spec B,\fq}
\end{tikzcd}
\]
由 \ref{it:1.2.4} 中的等同，此图变为
\[
\begin{tikzcd}
A \arrow[r,"\phi"] \arrow[d,"p_A"'] & B \arrow[d,"p_B"] \\
A_{f(\fq)} \arrow[r,"\phi'_{\fq}"'] & B_{\fq}
\end{tikzcd}
\]
其中竖直箭头 $p_A\colon A\to A_{f(\fq)}$ 与 $p_B\colon B\to B_{\fq}$ 是典范同态。于是
\[
\phi^{-1}(\fq)=\phi^{-1}p_B^{-1}(\fq B_{\fq})=p_A^{-1}(\phi'_{\fq})^{-1}(\fq B_{\fq})=p_A^{-1}(f(\fq)A_{f(\fq)})=f(\fq),
\]
其中第三个等号来自 $\phi'_{\fq}$ 是局部同态这一事实。因此有 $f(\fq)=\phi^{-1}(\fq)$。此外，最后一个图的交换性表明，$\phi'_{\fq}$ 与 $\phi$ 诱导的同态 $\phi_{\fq}\colon A_{\phi^{-1}(\fq)}\to B_{\fq}$ 一致。对 $\Spec A$ 的每个开子集 $V$ 与每个 $\fq\in f^{-1}(V)$，有交换图
\[
\begin{tikzcd}
\cO_{\Spec A}(V) \arrow[r,"f^\sharp"] \arrow[d] & \cO_{\Spec B}(f^{-1}(V)) \arrow[d] \\
\cO_{\Spec A,f(\fq)} \arrow[r,"f^\sharp_{\fq}"'] \arrow[d] & \cO_{\Spec B,\fq} \arrow[d] \\
A_{\phi^{-1}(\fq)} \arrow[r,"\phi_{\fq}"'] & B_{\fq}
\end{tikzcd}
\]
注意此图中竖直箭头的复合分别把截面送到它在 $\phi^{-1}(\fq)$ 处与在 $\fq$ 处的值。上面图的交换性表明：对 $\cO_{\Spec A}(V)$ 中的任意截面 $s\colon V\to\prod_{\fp\in\Spec A}A_{\fp}$，$\cO_{\Spec B}(f^{-1}(V))$ 中的截面 $f^\sharp(s)\colon f^{-1}(V)\to\prod_{\fq\in\Spec B}B_{\fq}$ 由
\[
f^\sharp(s)(\fq)=\phi_{\fq}(s(\phi^{-1}(\fq)))
\]
给出。因此 $(f,f^\sharp)$ 由 $\phi$ 诱导。
\end{proof}

\begin{tnote}
概形定义中的「局部」二字不可省。命题 \ref{it:1.2.5}(ii) 断言 $\Spec$ 之间的态射恰由环同态给出，其证明的关键一步是 $f^\sharp_{\fq}$ 为\emph{局部}同态：只有如此才能保证 $\phi'_{\fq}$ 把 $\fq B_{\fq}$ 的原像恰好拉回 $\fq A_{\fq}$，从而得到 $f(\fq)=\phi^{-1}(\fq)$。若把概形之间的态射放宽成任意环层空间的态射，$\Spec$ 就不再与环范畴反变等价，「仿射概形 $\leftrightarrow$ 环」这条对偶也就断了。[依据：原书 p.24–25 命题 1.2.5 的证明]
\end{tnote}

\begin{proposition}\label{it:1.2.6}
对任意 $f\in A$，有局部环层空间的典范同构
\[
(D(f),\cO_{\Spec A}|_{D(f)})\cong(\Spec A_f,\cO_{\Spec A_f}).
\]
\end{proposition}

\begin{proof}
典范同态 $A\to A_f$ 诱导态射 $(\varphi,\varphi^\sharp)\colon(\Spec A_f,\cO_{\Spec A_f})\to(\Spec A,\cO_{\Spec A})$。可以验证，映射 $\varphi\colon\Spec A_f\to\Spec A$ 是双射、连续、开映射，且像为 $D(f)$。故它诱导 $\Spec A_f$ 与 $D(f)$ 之间的同胚。然后可以证明，$\varphi^\sharp$ 诱导层的同构
\[
\cO_{\Spec A}|_{D(f)}\cong(\varphi_*\cO_{\Spec A_f})|_{D(f)},
\]
这用到 \ref{it:1.1.1} 以及对任意 $\fp\in D(f)$ 有 $A_{\fp}\cong(A_f)_{\fp}$。
\end{proof}

\begin{corollary}\label{it:1.2.7}
设 $(X,\cO_X)$ 是概形，$U$ 是 $X$ 的开子集。则 $(U,\cO_X|_U)$ 是概形。
\end{corollary}

我们称 $(U,\cO_X|_U)$ 为 $(X,\cO_X)$ 的\emph{开子概形}\idx{kaizigaixing}{开子概形}。

\begin{proof}
用仿射开子集 $U_i=\Spec A_i$（$i\in I$）覆盖 $X$。则 $U$ 可被 $U\cap U_i$（$i\in I$）覆盖。每个 $U\cap U_i$ 是 $U_i=\Spec A_i$ 的开子集，故可用 $D(f_{ij})$（$j\in J_i$）覆盖，其中 $f_{ij}\in A_i$。于是 $D(f_{ij})$（$i\in I$，$j\in J_i$）构成 $U$ 的开覆盖，且由 \ref{it:1.2.6}，每个 $(D(f_{ij}),\cO_X|_{D(f_{ij})})$ 是仿射概形。
\end{proof}

下面的命题是 \ref{it:1.2.5} 的推广：

\begin{proposition}\label{it:1.2.8}
设 $X$ 是概形，$A$ 是环。则从 $X$ 到 $\Spec A$ 的概形态射构成的集合，与从 $A$ 到 $\cO_X(X)$ 的环同态构成的集合之间，有一一对应。
\end{proposition}

\begin{proof}
定义映射
\[
\alpha\colon\Hom(X,\Spec A)\to\Hom(A,\cO_X(X))
\]
如下：对任意概形态射 $f\colon X\to\Spec A$，$\alpha(f)$ 为同态
\[
f^\sharp(\Spec A)\colon A\cong\cO_{\Spec A}(\Spec A)\to\cO_X(X).
\]
我们证明 $\alpha$ 是双射。用仿射开子概形 $U_i$（$i\in I$）覆盖 $X$，并用仿射开子概形 $U_{ijk}$（$k\in K_{ij}$）覆盖每个 $U_i\cap U_j$。概形态射 $f\colon X\to\Spec A$ 的集合 $\Hom(X,\Spec A)$ 可与满足 $f_i|_{U_{ijk}}=f_j|_{U_{ijk}}$（对一切 $i,j\in I$ 与 $k\in K_{ij}$）的元组 $(f_i)\in\prod_{i\in I}\Hom(U_i,\Spec A)$ 的集合等同。由命题 \ref{it:1.2.5}，此集合可与满足下图交换的元组 $(\phi_i)\in\prod_{i\in I}\Hom(A,\cO_X(U_i))$ 的集合等同：
\[
\begin{tikzcd}
A \arrow[r,"\phi_i"] \arrow[d,"\phi_j"'] & \cO_X(U_i) \arrow[d] \\
\cO_X(U_j) \arrow[r] & \cO_X(U_{ijk})
\end{tikzcd}
\]
即对一切 $i,j\in I$ 与 $k\in K_{ij}$ 交换。由于有正合序列
\[
0\to\cO_X(X)\to\prod_{i\in I}\cO_X(U_i)\to\prod_{i,j\in I,\,k\in K_{ij}}\cO_X(U_{ijk}),
\]
最后一个集合可与集合 $\Hom(A,\cO_X(X))$ 等同。
\end{proof}

回忆：\emph{分次环}\idx{fencihuan}{分次环}是环 $S$ 连同加法群 $S$ 的一个分解 $S=\bigoplus_{d=0}^{\infty}S_d$，即 $S$ 分解为交换群 $S_d$ 的直和，使得对任意非负整数 $d$ 与 $e$ 有 $S_dS_e\subset S_{d+e}$。$S_d$ 中的元素称为\emph{次数为 $d$ 的齐次元素}\idx{qiciyuansu}{齐次元素}。设 $1=\sum_ie_i$，其中 $e_i\in S_i$ 对一切 $i$ 成立。则对任意齐次元素 $x\in S$，由 $1\cdot x=x$ 得 $e_0x=x$，且对一切 $i\ge1$ 得 $e_ix=0$。于是 $e_0$ 是单位元，且与 $1$ 重合。所以 $1\in S_0$。$S$ 的理想 $\fa$ 称为\emph{齐次的}\idx{qicilixiang}{齐次理想}，若它满足以下等价条件之一：
\begin{enumerate}[label=(\roman*)]
  \item $\fa=\bigoplus_d(\fa\cap S_d)$。
  \item 若 $a\in\fa$ 且 $a=\sum_da_d$，其中 $a_d\in S_d$，则 $a_d\in\fa$。
  \item $\fa$ 作为 $S$ 的加法子群由齐次元素生成。
  \item $\fa$ 作为 $S$ 的理想由齐次元素生成。
\end{enumerate}

对 $S$ 的任意齐次理想 $\fa$，$S/\fa$ 是分次环，且典范同态 $S\to S/\fa$ 保持分次。

设 $\fa$ 是齐次理想。若对 $S$ 中任意齐次元素 $f$ 与 $g$，由 $fg\in\fa$ 可推出 $f\in\fa$ 或 $g\in\fa$，则 $\fa$ 是素理想。为看出这一点，设 $f$ 与 $g$ 是 $S$ 中的元素（不必齐次），满足 $fg\in\fa$ 但 $f,g\notin\fa$。设 $f=\sum_if_i$，$g=\sum_ig_i$，其中 $f_i$ 与 $g_i$ 对每个 $i$ 都是次数为 $i$ 的齐次元素。设 $m$ 是使 $f_m\notin\fa$ 的最小整数，$n$ 是使 $g_n\notin\fa$ 的最小整数。$fg$ 的次数 $m+n$ 分量为
\[
\sum_{i=0}^{m-1}f_ig_{m+n-i}+f_mg_n+\sum_{j=0}^{n-1}f_{m+n-j}g_j.
\]
由 $fg\in\fa$，此和属于 $\fa$。由 $m$ 与 $n$ 的取法，对 $i=0,\dots,m-1$ 有 $f_i\in\fa$，对 $j=0,\dots,n-1$ 有 $g_j\in\fa$。故必须有 $f_mg_n\in\fa$。由假设，$f_m\in\fa$ 或 $g_n\in\fa$，这与 $m$、$n$ 的取法矛盾。

我们把以下事实留给读者证明：齐次理想的和、积、交以及幂零根都是齐次理想。

令 $S_+=\bigoplus_{d=1}^{\infty}S_d$，令 $\Proj S$ 为 $S$ 的不包含 $S_+$ 的齐次素理想构成的集合。对 $S$ 的任意齐次理想 $\fa$，定义
\[
V_+(\fa)=\{\fp\in\Proj S\mid \fa\subset\fp\}.
\]
下面的命题留给读者证明：
\begin{proposition}\label{it:1.2.9}
\begin{enumerate}[label=(\roman*)]
  \item $V_+(0)=\Proj S$，$V_+(S)=\emptyset$。
  \item 对 $S$ 的任意一族齐次理想 $\fa_i$（$i\in I$），有 $\bigcap_{i\in I}V_+(\fa_i)=V_+\bigl(\sum_{i\in I}\fa_i\bigr)$。
  \item 对 $S$ 的任意齐次理想 $\fa$ 与 $\fb$，有 $V_+(\fa)\cup V_+(\fb)=V_+(\fa\fb)=V_+(\fa\cap\fb)$。
\end{enumerate}
\end{proposition}

这个命题表明，$\Proj S$ 是拓扑空间，其闭集恰为齐次理想 $\fa$ 对应的 $V_+(\fa)$。这个拓扑称为 $\Proj S$ 上的 \emph{Zariski 拓扑}。

设 $T$ 是 $S$ 的\emph{乘性子集}，即 $T$ 是包含 $1$ 且对乘法封闭的 $S$ 的子集。则
\[
\Bigl\{\frac{a}{t}\in T^{-1}S\Bigm| a\in S,\ t\in T,\ a \text{ 与 } t \text{ 是同次数的齐次元素}\Bigr\}
\]
是 $T^{-1}S$ 的子环。当 $T=S-\fp$ 而 $\fp$ 是齐次素理想时，把这个环记为 $S_{(\fp)}$。当 $T=\{1,f,f^2,\dots\}$ 而 $f\in S$ 是齐次元素时，把这个环记为 $S_{(f)}$。

在 $\Proj S$ 上定义环层 $\cO_{\Proj S}$ 如下：对 $\Proj S$ 的每个开子集 $U$，$\cO_{\Proj S}(U)$ 由满足以下两个条件的映射 $s\colon U\to\prod_{\fp\in\Proj S}S_{(\fp)}$ 组成：
\begin{enumerate}[label=(\alph*)]
  \item 对每个 $\fp\in U$，有 $s(\fp)\in S_{(\fp)}$。
  \item 对每个 $\fp\in U$，存在 $\fp$ 的包含在 $U$ 中的邻域 $U_{\fp}$ 及同次数的齐次元素 $a,f\in S$，使得对每个 $\fq\in U_{\fp}$，有 $f\notin\fq$ 且 $s(\fq)=\frac{a}{f}\in S_{(\fq)}$。
\end{enumerate}

对开子集的每个包含 $V\subset U$，定义 $\cO_{\Proj S}(U)\to\cO_{\Proj S}(V)$ 为映射的限制。若无混淆，常把 $\cO_{\Proj S}$ 记为 $\cO$。

\begin{proposition}\label{it:1.2.10}
\begin{enumerate}[label=(\roman*)]
  \item 对每个 $\fp\in\Proj S$，有典范同构 $\cO_{\fp}\cong S_{(\fp)}$。
  \item 对每个正次数的齐次元素 $f\in S_+$，令
  \[
  D_+(f)=\Proj S-V_+((f))=\{\fp\in\Proj S\mid f\notin\fp\}.
  \]
  则 $D_+(f)$ 是 $\Proj S$ 中的开集。这类开集构成 $\Proj S$ 的拓扑的一组基。此外，有局部环层空间的同构
  \[
  (D_+(f),\cO_{\Proj S}|_{D_+(f)})\cong(\Spec S_{(f)},\cO_{\Spec S_{(f)}}).
  \]
  特别地，$(\Proj S,\cO_{\Proj S})$ 是概形。
\end{enumerate}
\end{proposition}

\begin{proof}
我们证明 (ii)，把 (i) 留给读者。作为闭集 $V_+((f))$ 的补集，$D_+(f)$ 是开集。对 $S$ 的每个齐次理想 $\fa$ 与开子集 $\Proj S-V_+(\fa)$ 中的每个点 $\fp$，有 $\fa\not\subset\fp$、$S_+\not\subset\fp$，从而 $\fa S_+\not\subset\fp$。取 $\fa S_+$ 中不属于 $\fp$ 的齐次元素 $f$。则 $f$ 有正次数，且 $\fp\in D_+(f)\subset\Proj S-V_+(\fa)$。故形如 $D_+(f)$ 的开子集构成一组基。

考虑映射
\[
\varphi\colon D_+(f)\to\Spec S_{(f)},\qquad \fp\mapsto \fp S_f\cap S_{(f)}.
\]
我们断言 $\varphi$ 是双射。事实上，设 $\fp_1,\fp_2\in D_+(f)$ 且 $\varphi(\fp_1)=\varphi(\fp_2)$。则对任意齐次元素 $x\in\fp_1$，有
\[
\frac{x^{\deg f}}{f^{\deg x}}\in\fp_1S_f\cap S_{(f)}=\fp_2S_f\cap S_{(f)}.
\]
（这里用到 $\deg f>0$。）故 $x^{\deg f}\in\fp_2$，从而 $x\in\fp_2$。因此 $\fp_1\subset\fp_2$。同理 $\fp_2\subset\fp_1$。这证明了 $\varphi$ 的单射性。对任意 $\fq\in\Spec S_{(f)}$，设 $\fp$ 是由那些满足「对某个 $k$ 有 $\frac{a}{f^k}\in\fq$」的元素 $a\in S$ 生成的理想的幂零根。（于是 $a$ 必是次数为 $k\deg f$ 的齐次元素。）注意 $\fp$ 是齐次的。设 $a_1$ 与 $a_2$ 是齐次元素且 $a_1a_2\in\fp$。则存在幂 $(a_1a_2)^n$ 落在由 $\fq$ 中元素的分子生成的理想中。注意 $(a_1a_2)^{n\deg f}$ 落在同一个理想中。这说明 $\frac{(a_1a_2)^{n\deg f}}{f^{n(\deg a_1+\deg a_2)}}$ 属于 $\fq$。由 $\fq$ 是素理想，有 $\frac{a_1^{n\deg f}}{f^{n\deg a_1}}\in\fq$ 或 $\frac{a_2^{n\deg f}}{f^{n\deg a_2}}\in\fq$。故 $a_1\in\fp$ 或 $a_2\in\fp$。因此 $\fp$ 是素理想。显然 $f\notin\fp$，故 $\fp\in D_+(f)$。可以验证 $\varphi(\fp)=\fq$。所以 $\varphi$ 是满射。

对 $S$ 的任意齐次理想 $\fa$，令 $\fb=\fa S_f\cap S_{(f)}$。可以验证 $\varphi(D_+(f)\cap V_+(\fa))=V(\fb)$。另一方面，对 $S_{(f)}$ 的任意理想 $\fb$，定义 $\fa$ 为由 $\fb$ 中元素的分子生成的 $S$ 的理想。则 $\fb=\fa S_f\cap S_{(f)}$。这说明 $\varphi$ 建立了 $D_+(f)$ 的闭子集族与 $\Spec S_{(f)}$ 的闭子集族之间的一一对应。所以 $\varphi$ 是同胚。

对每个 $\fp\in D_+(f)$，有显然的同态
\[
(S_{(f)})_{\varphi(\fp)}\to S_{(\fp)}.
\]
给定元素 $\frac{a}{t}\in S_{(\fp)}$，有 $t\notin\fp$ 且 $\deg a=\deg t$。于是 $\frac{t^{\deg f}}{f^{\deg t}}\in S_{(f)}-\varphi(\fp)$，且 $\frac{t^{\deg f-1}a}{f^{\deg t}}\in S_{(f)}$。容易验证
\[
\frac{\frac{t^{\deg f-1}a}{f^{\deg t}}}{\frac{t^{\deg f}}{f^{\deg t}}}
\]
落在 $(S_{(f)})_{\varphi(\fp)}$ 中，并被送到 $S_{(\fp)}$ 中的 $\frac{a}{t}$。这证明了 $(S_{(f)})_{\varphi(\fp)}\to S_{(\fp)}$ 是满射。不难证明它也是单射。利用这个同构，可以构造层的同构 $\varphi^\sharp\colon\cO_{\Spec S_{(f)}}\cong\varphi_*(\cO_{\Proj S}|_{D_+(f)})$。
\end{proof}

\begin{tnote}
命题 \ref{it:1.2.10}(ii) 是全书反复使用的「射影概形局部仿射」的来源：$\Proj S$ 被 $D_+(f)$ 覆盖（$f$ 取遍正次数的齐次元素），而每个 $D_+(f)$ 同构于仿射概形 $\Spec S_{(f)}$。注意分次局部化 $S_{(f)}$ 只取「分子与分母同次数」的分式，所以它不带 $\bZ$-分次，是一个普通环；这正是 (ii) 能把 $\Proj$ 的问题化归到 $\Spec$ 的原因。[依据：原书 p.27–29 命题 1.2.10 及其证明]
\end{tnote}

\begin{exercises}
\item \label{ex:1.2.1} 设 $A$ 是环。对 $\Spec A$ 的每个开子集 $U$，令 $S_U$ 为乘性子集 $S_U=\bigcap_{\fp\in U}(A-\fp)$，并令 $\cP(U)=S_U^{-1}A$。对开子集的每个包含 $V\subset U$，有 $S_U\subset S_V$，从而有典范同态 $\cP(U)\to\cP(V)$。这使得 $\cP$ 成为 $\Spec A$ 上环的前层。证明 $\cO_{\Spec A}\cong\cP^+$。
\item \label{ex:1.2.2} 设 $S$ 是环 $A$ 的乘性子集。证明典范态射 $\Spec S^{-1}A\to\Spec A$ 在底拓扑空间上诱导嵌入。
\item \label{ex:1.2.3} 设 $x$ 是概形中的点，$k(x)=\cO_{X,x}/\fm_x$ 是 $x$ 处的剩余域。构造自然态射 $i\colon\Spec k(x)\to X$，使其像为 $x$，且 $i^\sharp$ 诱导的同态 $\cO_{X,x}\to k(x)$ 是典范同态。
\item \label{ex:1.2.4} 设 $S$ 是分次环。
\begin{enumerate}[label=(\roman*)]
  \item 设 $\fp$ 是 $S$ 的素理想，令
  \[
  \fp'=\bigoplus_d(\fp\cap S_d).
  \]
  证明 $\fp'$ 是 $S$ 的齐次素理想。
  \item 设 $\fa$ 是 $S$ 的齐次理想。证明 $\sqrt{\fa}$ 是包含 $\fa$ 的齐次素理想的交。
\end{enumerate}
\item \label{ex:1.2.5} 设 $\phi\colon S\to T$ 是分次环的同态。设存在整数 $m\ge1$，使得对一切 $d$ 有 $\phi(S_d)\subset T_{md}$。设 $I$ 是由 $\phi(S_+)$ 生成的 $T$ 的齐次理想，$U=\Proj T-V(I)$。构造概形态射 $f\colon U\to\Proj S$，使得对任意 $\fq\in U$ 有 $f(\fq)=\phi^{-1}(\fq)$，且 $f^\sharp_{\fq}\colon\cO_{\Proj S,f(\fq)}\to\cO_{\Proj T,\fq}$ 可与同态 $\phi_{\fq}\colon S_{(\phi^{-1}(\fq))}\to T_{(\fq)}$ 等同。
\item \label{ex:1.2.6} 设 $A$ 是环，$I$ 是 $A$ 的理想，$S=\bigoplus_{d=0}^{\infty}I^d$。则 $S$ 是分次环。我们称 $\Proj S$ 为 $\Spec A$ 沿理想 $I$ 的\emph{爆破}\idx{baopo}{爆破}。证明包含映射 $A=S_0\hookrightarrow S$ 诱导概形态射 $f\colon\Proj S\to\Spec A$，使得在开子集 $U=\Spec A-V(I)$ 上，$f$ 诱导同构 $f^{-1}(U)\xrightarrow{\sim}U$。
\item \label{ex:1.2.7} 设 $A=R[x_1,\dots,x_n]$ 是多项式环，$I=(x_1,\dots,x_n)$ 是由 $x_1,\dots,x_n$ 生成的 $A$ 的理想，$S=\bigoplus_{d=0}^{\infty}I^d$，$T=A[y_1,\dots,y_n]$ 是分次环，使得 $T_d$ 由系数在 $A$ 中的 $y_1,\dots,y_n$ 的 $d$ 次齐次多项式组成，$J$ 是由 $x_iy_j-x_jy_i$（$i,j=1,\dots,n$）生成的 $T$ 的齐次理想。考虑分次环的满同态 $\phi\colon T\to S$，使得对任意 $a\in A$ 有 $\phi(a)=a\in S_0$，且 $\phi(y_i)=x_i\in S_1$。证明 $\phi$ 诱导概形的同构 $\Proj S\xrightarrow{\sim}\Proj(A[y_1,\dots,y_n]/J)$。
\end{exercises}
